\section{Quantitative half-line approximation and frequency tails}

This section keeps the nonprincipal primitive Dirichlet sector
and the original finite prolate family of AC-T38--T39.
All constants may depend on this fixed character.
The local factorization and frequency bump estimates used below
are those of AC-T40.

\begin{lemma}[AC-T41-ratio: a zero-gap-free shifted-line estimate]
\label{AC-T41-ratio}
Assume \(\mathsf A\), and write
\[
 M_\delta(t)=
  \frac{L(1/2+it,\chi)}{L(1/2+\delta+it,\chi)},
                         \qquad 0<\delta\leq1/4.
 \tag{TT.1}
\]
Then
\[
 |M_\delta(t)|\leq C_\chi(2+|t|)^{\delta/2},
 \qquad
 \|M_\delta-1\|_{L^2([-K,K])}
       \leq C_\chi K\log(2+K)\sqrt\delta
 \tag{TT.2}
\]
for \(K\geq2\), \(\delta\leq1/K\).
There is also the uniform crude bound
\[
 \frac1{|L(\sigma+it,\chi)|}
       \leq \{C_\chi(2+|t|)\}^{C_\chi/\delta},
                  \qquad \sigma\geq1/2+\delta.
 \tag{TT.3}
\]
No zero-separation or derivative-at-a-zero estimate is assumed.
\end{lemma}
\begin{proof}
The functional equation and conjugation imply that
\[
        F(t)=e^{-i\arg(W_\chi)/2}
                     \Lambda_\chi(1/2+it)
 \]
is real on the real axis.  It is a real entire function of
order one.  Under \(\mathsf A\) all of its zeros are real.
The usual completed function is entire in this nonprincipal
sector, so no polar factor is being suppressed.
Hadamard factorization, with each zero repeated to its full
multiplicity, gives
\[
 F(z)=e^{A+Bz}z^{m_0}
       \prod_{\gamma\ne0}(1-z/\gamma)e^{z/\gamma},
       \qquad B\in\mathbb R,\qquad
                   \sum_{\gamma\ne0}\gamma^{-2}<\infty.
 \tag{TT.4}
\]
The constant phase in \(e^A\) does not affect ratios.
The reality of \(B\) follows by conjugating this factorization.
For exact background locators, the Hadamard theorem is
Elkies, \emph{Functions of finite order}, Theorem (1),
pp.~1--4, at
\url{https://abel.math.harvard.edu/~elkies/M229.22/prod.pdf};
the nonprincipal completed Dirichlet product is equation (1),
p.~1, of
\url{https://people.math.harvard.edu/~elkies/M229.19/pnt_q.pdf}.
These are the standard analytic inputs to (TT.4); the
shifted-line estimates below are proved here.
Thus every factor of \(F(t)/F(t-i\delta)\) has modulus at
most one.  This proves
\[
                 |F(t)/F(t-i\delta)|\leq1.
 \tag{TT.5}
\]

Here is a quantitative comparison with one on a compact
interval.  For a zero \(\gamma\),
\[
 \int_{\mathbb R}
   \left|\frac{t-\gamma}{t-\gamma-i\delta}-1\right|^2dt
                         =\pi\delta .
 \tag{TT.6}
\]
For \(|\gamma|\leq2K+2\), telescope the corresponding
finite product.  All other factors in each telescoping
term have modulus at most one.
There are \(O_\chi(K\log(2+K))\) such zeros.
This count, and the two partial-summation consequences
used below, follow also by covering the critical line
with the fixed disks in (FB.6) of AC-T40; the bounded
height portion contains only finitely many zeros.
The genus-one phase \(e^{i\delta/\gamma}\) contributes at
most
\(C_\chi\delta\sqrt K\log^2(2+K)\);
the finitely many zeros near zero and the possible central
factor are absorbed in \(C_\chi\).
For \(|\gamma|>2K+2\), expand the logarithm of the
normalized ratio factor.  Its linear \(1/\gamma\) term
cancels the genus-one phase, and uniformly for \(|t|\leq K\)
the remainder is at most
\[
                     C\delta(K+1)/\gamma^2.
 \]
The zero count gives
\(\sum_{|\gamma|>2K+2}\gamma^{-2}
       =O_\chi(\log(2+K)/K)\).
The far product therefore differs from one by
\(O_\chi(\delta\log(2+K))\) when \(\delta\leq1/K\).
The exponential \(e^{iB\delta}\) is harmless.
Equation (TT.6), the triangle inequality in \(L^2\), and
these two phase estimates prove
\[
 \left\|\frac{F(t)}{F(t-i\delta)}-1\right\|_{L^2([-K,K])}
               \leq C_\chi K\log(2+K)\sqrt\delta .
 \]
The gamma and conductor factors which change this completed
ratio into \(M_\delta\) have modulus at most
\(C_\chi(2+|t|)^{\delta/2}\), by Stirling's formula and
the absence of gamma poles on the stated strip.
Their difference from one on \([-K,K]\) is
\(O_\chi(\delta\log(2+K))\).
This proves both estimates in (TT.2).

For (TT.3), use the analytic logarithm of \(L\) on the
zero-free half-plane \(\operatorname{Re}s>1/2\).
Center a disk at \(2+it\), with radius
\(3/2-\delta/2\).
The polynomial strip bound bounds \(\log|L|\) above
there by \(C_\chi\log(2+|t|)\).
Subtract it from a sufficiently large constant multiple
of that logarithm to obtain a nonnegative harmonic
function.  Its value at the center is
\(O_\chi(\log(2+|t|))\), since the Euler product bounds
\(|L(2+it)|\) above and below.
Harnack's inequality, at any point of the horizontal
segment from \(1/2+\delta+it\) to \(2+it\), costs at
most \(C/\delta\).
It follows that
\(-\log|L|\leq C_\chi\delta^{-1}\log(2+|t|)\)
on this segment.  The Euler product handles
\(\operatorname{Re}s\geq2\).
Exponentiation proves (TT.3).
\end{proof}

\begin{lemma}[AC-T41-halfline: quantitative approximation of a half-line column operator]
\label{AC-T41-halfline}
Assume \(\mathsf A\).
There are constants \(B_\chi,C_{\chi,m}\) with the following
property.  Let \(\mathcal A\) be any Hilbert space and let
\(F:\mathcal A\to L^2((-\infty,0])\) satisfy
\[
 \|F\|_2+\|F\|_{2,\mathcal A\to H^{1/4}(\mathbb R)}
             +\|tF\|_2\leq B,\qquad B\geq B_\chi ,
 \tag{TT.7}
\]
where the half-line outputs are extended by zero.
There is an actual raw parity column operator
\(\Psi:\mathcal A\to X_m\), with every column smoothly
supported in a common compact subinterval of
\(\{0<|u|<1\}\), such that
\[
 \|E_\chi\Psi-F\|_2\leq1/100,\qquad
 \|\Psi\|_{2,\mathcal A\to X_m}
                      \leq \exp(C_{\chi,m}B^{102}).
 \tag{TT.8}
\]
Here \(X_m\) is the spatial Sobolev space from AC-S29.
Taking \(m=12\), for the actual time-endpoint dilation, or
for the dual-character Fourier-endpoint construction,
the approximation by the selected finite images has error
at most \(1/50\) as soon as
\[
                    \log\Lambda\geq C_{\chi}B^{102}.
 \tag{TT.9}
\]
Changing the two fixed error tolerances only changes constants.
\end{lemma}
\begin{proof}
All estimates below are Hilbert--Schmidt column estimates;
they follow from the corresponding scalar inequalities by
summing over an orthonormal input basis.
Fix smooth compact cutoffs and a nonnegative unit-mass
mollifier whose \(n\)-th derivative is bounded by
\(\exp(C(n+1)\log(n+2))\).
Such fixed functions can be obtained from
\(\exp[-1/(1-x^2)]1_{|x|<1}\).
For example Cauchy's estimate on a complex disk of radius
a fixed multiple of the square of the distance to the
endpoint gives \(n!C^n d^{-2n}e^{-c/d}\);
optimizing in \(d\) proves the displayed bound.
This elementary bound, rather than an unspecified
smoothness constant, is used below.

Take
\[
       M=C_0B,\qquad \eta=(C_0B)^{-4},
                         K=B^{10},\qquad\delta=B^{-50},
 \tag{TT.10}
\]
where \(C_0\), and then \(B_\chi\), are chosen large enough.
Convolve \(F\) with a probability mollifier supported in
\([-3\eta,-2\eta]\).
Fourier translation estimates in \(H^{1/4}\) bound its
distance from \(F\) by \(C\eta^{1/4}B\).
Its first log moment remains bounded by \(2B\).
Multiply by a smooth cutoff equal to one on \([-M,0]\)
and zero below \(-2M\).
The first-moment bound makes this further error at most
\(2B/M\).
The resulting column operator \(F_0\) has output support
in \([-2M-3\eta,-2\eta]\) and satisfies
\[
 \begin{split}
 \|F_0-F\|_2&\leq1/1000,\\
 \|F_0\|_{2,H^1}&\leq C B^5,\\
 \|F_0\|_{2,H^n}
      &\leq \exp\!\bigl(C(n+1)
                       \{\log(n+2)+\log(B+2)\}\bigr).
 \end{split}
 \tag{TT.11}
\]
Increasing the fixed \(C_0\) gives the first numerical
tolerance.  Leibniz's rule, the fixed cutoff derivative
bound, and the scale \(\eta\) give the last two estimates.

Let \(\widehat F_0(z)\) denote its entire Fourier row
operator and define
\[
        \widehat h_\delta(t)
          =\frac{\widehat F_0(t-i\delta)}
                         {L(1/2+\delta+it,\chi)} .
 \tag{TT.12}
\]
This defines an ordinary \(L^2\) column operator.
Indeed (TT.3) gives a finite polynomial growth exponent
\(C_\chi/\delta\), and (TT.11) gives arbitrarily many
derivatives, with explicit bounds.
More precisely, on each line
\(\operatorname{Im}z=v\), \(0\leq v\leq\delta/2\),
use \(n=\lceil 2C_\chi/\delta\rceil+j+4\) in (TT.11).
Plancherel and (TT.3) give
\[
 \|\widehat h_\delta(\,\cdot+i\delta/2)\|
                _{2,(1+t^2)^{j/2}}
             +\|h_\delta\|_{2,H^j}
                 \leq \exp(C_{\chi,j}B^{51})
                  \qquad(j\text{ fixed}).
 \tag{TT.13}
\]
Here \(n=O_{\chi,j}(B^{50})\) and
\(n\log(nB)=O_{\chi,j}(B^{50}\log B)\leq C_{\chi,j}B^{51}\).

The support statement must also be justified.
The quotient in (TT.12) is analytic on
\(\operatorname{Im}z<\delta\).
In the lower half-plane its inverse \(L\)-factor is bounded
by (TT.3) on bounded horizontal strips, and by the absolutely
convergent Euler product sufficiently far below.
The support gap of \(F_0\) supplies an exponential factor
\(e^{2\eta\operatorname{Im}z}\), up to polynomial factors
from finitely many derivatives.
Shifting the inverse Fourier contour down to
\(\operatorname{Im}z=-V\), and then letting \(V\to\infty\),
therefore proves
\[
                         h_\delta(t)=0\quad(t>-\eta).
 \tag{TT.14}
\]
Equivalently, after a translation by \(\eta\), this is
the ordinary half-plane Hardy support argument.
The integrals on the horizontal contours are absolutely
convergent by the same derivative estimates; on the remote
lower contours the Euler product bound applies.
More explicitly, at \(z=x-iV\) the numerator and its
fixed derivative bounds have the factor
\(\operatorname{poly}(V)e^{-2\eta V}\).
Inverse Fourier transformation contributes \(e^{-Vt}\).
For \(t>-\eta\) their product tends to zero at least
as a polynomial times \(e^{-\eta V}\).
The vertical contour pieces vanish by taking a fixed
sufficiently large number of frequency integrations by
parts before sending their real endpoints to infinity.
The positive contour shift in (TT.13), separately, gives
\[
       \|e^{-\delta t/2}h_\delta(t)\|_{2,H^j}
                           \leq\exp(C_{\chi,j}B^{51}).
 \tag{TT.15}
\]
Thus support and negative-tail decay have different,
explicit justifications.

We next check the approximation in the actual theta norm.
The Mellin identity gives
\[
 \mathcal F_+(E_\chi U^{-1}h_\delta)
               =M_\delta(t)\widehat F_0(t-i\delta).
 \tag{TT.16}
\]
Damping \(F_0\) by \(e^{\delta t}\) changes it in
Hilbert--Schmidt norm by at most \(C\delta MB\).
On \([-K,K]\), the Fourier row of this damped operator
has norm at most \(C\sqrt M B\), by Cauchy--Schwarz on
its support interval.
The second estimate (TT.2) therefore bounds the
low-frequency Hilbert--Schmidt error by
\[
              C_\chi \sqrt M B\,K\log(2+K)\sqrt\delta.
 \tag{TT.17}
\]
On the complement, the first estimate (TT.2) and the
\(H^1\) bound in (TT.11) give an error at most
\[
                        C_\chi K^{-7/8}B^5 .
 \tag{TT.18}
\]
The three quantities are respectively
\(O_\chi(B^{-48})\),
\(O_\chi(B^{-13.5}\log B)\), and
\(O_\chi(B^{-3.75})\), after (TT.10).
For \(B_\chi\) sufficiently large they are smaller than
the fixed tolerances in (TT.8).
This also verifies directly that the vector in (TT.16)
belongs to the closed theta domain.

Finally choose \(H=C_{\chi,m}B^{102}\) and smoothly cut
\(h_\delta\) off below \(-H\), using a transition of fixed
logarithmic width one.  Equations (TT.15) and (TU.5)
bound the discarded theta graph norm by
\[
       C_\chi e^{-\delta H/4}
                           \exp(C_{\chi,m}B^{51}),
 \]
which is smaller than the required tolerance.
The remaining log profile is supported in
\([-H-1,-\eta]\).
On returning to the raw coordinate,
each ordinary derivative of order at most \(m\) costs
at most a constant times \(e^{(m+1)(H+1)}\), together
with the log derivative bounds already in (TT.13).
The resulting raw parity column operator \(\Psi\)
therefore satisfies (TT.8).
It is smooth, is zero near raw zero, and has the required
compact support strictly inside the time interval.

For \(m=12\), AC-S29 gives the actual graph error
\[
  \|E_\chi(1-\Pi_\Lambda)D_\Lambda^{\rm raw}\Psi\|_2
       \leq C_\chi c^{-27/16}\exp(C_\chi B^{102}),
 \tag{TT.19}
\]
where \(D_\Lambda^{\rm raw}\) is the raw unitary dilation.
This proves (TT.9), after adjusting its constant.
For the lower endpoint, apply the same construction to
the reflected target operator in the dual character,
then Fourier transform the raw profiles back and dilate
their Fourier support to radius \(R=\Lambda/q\).
Their time tails have the same spatial-Sobolev estimate
as in the second profile assertion of AC-T38.
The identical interpolation gives (TT.19).
Primitive Poisson gives the target endpoint
\(1/(qR)=1/\Lambda\) exactly.
\end{proof}

\begin{theorem}[AC-T41: polynomial local mass and the full fixed-sector trace]
\label{AC-T41}
Assume \(\mathsf A\), and retain a fixed nonprincipal primitive
Dirichlet character \(\chi\).
There are \(A_\chi,C_\chi<\infty\) such that the actual
positive defect measures satisfy
\[
        \mu_\Lambda([T,T+1])
            \leq C_\chi(2+|T|)^{A_\chi}
                       \qquad(\Lambda\geq2,\ T\in\mathbb R).
 \tag{TT.20}
\]
Consequently their convergence in AC-T39 holds against
every Schwartz function on the frequency line.
For every original \(h\in C_c^\infty(\mathbb R_+)\),
with the original multiplicative measure and Fourier
conventions, one obtains
\[
 \operatorname{Tr}(Q_\Lambda V(h))
       =2\log\Lambda\,h(1)
           -\sum_\gamma m_{\chi,\gamma}\widehat h_+(\gamma)
              +o(1).
 \tag{TT.21}
\]
The sum retains every zero multiplicity.
This is the actual fixed-sector finite-prolate formula.
Its scope is nonprincipal primitive Dirichlet sectors;
the principal polar corrections and the exact full
adelic cutoff dictionary are not added by this theorem.
\end{theorem}
\begin{proof}
We give the uniformity argument, since a merely qualitative
version of the local upper proof would not imply (TT.20).
For \(|T|\) large apply AC-T40 with center \(T+1/2\),
and denote its adaptive flat-top bump by \(a_T\).
Thus \(a_T=1\) on \([T,T+1]\).
It has fixed compact frequency support around that interval,
\(0\leq a_T\leq1\), and uses \(N_T=O_\chi(\log(2+|T|))\)
local zero factors, with their multiplicities.
Let \(P_{Z,L,T}\) be the window projection onto those
finite zero jets and let \(W_{L,T}=S_L-P_{Z,L,T}\).
No lower bound on this jet Gram matrix is needed here:
\[
        \|W_{L,T}\|\leq1,\qquad
                    \operatorname{rank}P_{Z,L,T}\leq N_T.
 \tag{TT.22}
\]

Fix, for example, the derivative and moment budgets
\(J=k=20\) in AC-T40; these exceed all fixed budgets
used below.
Use the finite-pole inverse \(T_{L,T}\) in the proof of
AC-T39, now with the stable local Blaschke factorization
and adaptive bump of AC-T40.
Its estimates give constants \(C,d\), depending only on
the fixed character and fixed derivative budgets, such that
\[
 \|T_{L,T}\|_2
       \leq C(2+|T|)^d(1+L)^{N_T+d}.
 \tag{TT.23}
\]
Both upper raw boundary kernels, including the Fourier-dual
one, have the fixed weighted output derivative norms needed
in (TU.15)--(TU.16), bounded by \(C(2+|T|)^d\).
This is exactly the gap-free positive-tail estimate in
AC-T40; the negative polynomial tail is canceled by
the finite zero-jet conditions in (TT.22).

One can keep the jet correction throughout this argument
as a right multiplication by \(W_{L,T}\).
For example the uncorrected translated upper kernel is
the fixed-quadrant kernel in (TU.15), restricted to
\([-2L,0]\) in its input, and then composed with the
translated zero-jet complement.
Both right multipliers are contractions.
Thus all its weighted output Hilbert--Schmidt norms
remain bounded by the uncorrected kernel norms.
This avoids the separation-dependent convergence of
individual jet Gram coefficients.

After applying the closed theta map, the first target
log derivative and the first target log moment of these
boundary operators are still bounded by a fixed power
of \(2+|T|\).  To see the moment assertion explicitly,
Fourier differentiation gives
\[
       \partial_\tau(L_\chi\,\widehat\psi)
          =L_\chi'\widehat\psi
                                 +L_\chi\partial_\tau\widehat\psi.
 \]
The unconditional polynomial bounds for \(L_\chi,L_\chi'\)
and the fixed derivative and moment budget in AC-T40
control both terms by Plancherel.
The corresponding statement for the derivative is even
more immediate.
Sharp target half-line or interval restriction has
uniform \(H^{1/4}\) control by the two-endpoint estimate
(TB.103), applied to these whole-line \(H^1\) outputs.
Restriction does not increase the weighted \(L^2\) moment.
Hence every target boundary term in (TU.17), after its
appropriate endpoint translation, satisfies (TT.7) with
\[
                         B_T\leq C_\chi(2+|T|)^d.
 \tag{TT.24}
\]
The same holds for the two input boundary terms
\(S_LC_{a_T}(1-S_L)\).
Their fixed-quadrant convolution kernels use only a
fixed number of derivatives of the adaptive bump,
whose bounds are polynomial in \(2+|T|\).

For completeness the cross errors are uniform as well.
The proof of (TU.16), with any fixed sufficiently large
ordinary derivative budget, now bounds them by
\[
                         C(2+|T|)^d c^{-2}.
 \tag{TT.25}
\]
The bulk approximation (TU.2) and (TT.23) bound its
graph error, after increasing a fixed Fourier-gap
derivative budget, by
\[
          C(2+|T|)^d(1+L)^{N_T+d}e^{-L}.
 \tag{TT.26}
\]
There is a fixed exponent \(D\) for which both errors
are less than \(1/100\) whenever
\[
                         L\geq C_\chi(2+|T|)^D.
 \tag{TT.27}
\]
Indeed \(N_T=O(\log(2+|T|))\), and the logarithm of
the polynomial-in-\(L\) factor in (TT.26) is
\(O(\log(2+|T|)\log(1+L))\), which is absorbed by \(L/2\)
uniformly beyond such a threshold.

Apply the quantitative half-line lemma to each of the
finitely many boundary terms, with (TT.24), using the
dual character for the lower ones.
Its condition \(\log\Lambda\geq C_\chi B_T^{102}\)
is another condition of the form (TT.27), after enlarging
the single fixed exponent \(D\).
The exact residual decomposition (TU.17), including
(TU.16a), therefore proves
\[
       \|(1-Q_\Lambda)S_LC_{a_T}
                                (1-P_{Z,L,T})\|_2\leq1
               \quad\text{whenever (TT.27) holds}.
 \tag{TT.28}
\]
The input splitting is orthogonal.
By (TT.22) and \(\|C_{a_T}\|\leq1\), its zero-jet
summand has squared Hilbert--Schmidt norm at most \(N_T\).
The exact distance formula (TU.22) yields
\[
       \mu_\Lambda([T,T+1])
          \leq \int|a_T|^2\,d\mu_\Lambda
          \leq N_T+1
                   \quad\text{under (TT.27)}.
 \tag{TT.29}
\]

In the complementary range, the elementary density
bound \(0\leq d\mu_\Lambda/d\tau\leq L/\pi\) gives
\[
       \mu_\Lambda([T,T+1])
                       \leq L/\pi
                       \leq C_\chi(2+|T|)^D.
 \tag{TT.30}
\]
This is why a threshold polynomial in \(\log\Lambda\),
as expressed in (TT.27), suffices.
For bounded \(T\), AC-T39 on one fixed larger compact
frequency interval gives a uniform bound for all
sufficiently large \(\Lambda\); the same elementary
density bound handles the remaining bounded cutoffs.
This proves (TT.20).

Let \(\varphi\) be Schwartz.  On integer unit intervals
the right side of (TT.20) is polynomial, whereas
\(\sup_{[n,n+1]}|\varphi|\) decreases faster than every
power.  Summing these bounds shows
\[
 \lim_{K\to\infty}\sup_{\Lambda\geq2}
             \int_{|\tau|>K}|\varphi(\tau)|\,d\mu_\Lambda(\tau)=0.
 \tag{TT.31}
\]
The zero-counting measure has the same integrability,
by its standard polynomial zero count.
Compact-frequency convergence from AC-T39 and
(TT.31) prove convergence against every Schwartz
function.

For an original compact smooth norm test \(h\),
\(\widehat h_+\) is Schwartz.
The identity defining \(\mu_\Lambda\) gives
\[
 \int\widehat h_+\,d\mu_\Lambda
       =2Lh(1)-
             \operatorname{Tr}(S_\Lambda Q_\Lambda
                                     S_\Lambda V(h)).
 \]
Finally AC-T30 replaces the compressed contraction
by the actual \(Q_\Lambda\) in this tested trace,
with error tending to zero for this original test.
The established Schwartz convergence proves (TT.21).
\end{proof}

\begin{remark}[Trace and hypothesis scope]
The proof uses \(\mathsf A\) for the real-zero canonical
product and for the zero-free half-plane in (TT.3).
It makes no independent claim about a formal verification
of \(\mathsf A\).
All approximants are actual Hilbert--Schmidt column
operators, and all final projections are the original
finite arithmetic theta-image projections.
No completed range replaces that finite projection.
The result does not claim RH for a non-Dirichlet Hecke
or Artin sector, and it does not suppress the two polar
corrections needed in the principal rational sector.
\end{remark}
